\documentclass[letterpaper,twocolumn,10pt]{article}
\usepackage{usenix-2020-09}
\usepackage{booktabs}
\usepackage{graphicx}
\usepackage{url}
\usepackage{xspace}
\usepackage{amsmath}
\usepackage{enumitem}

\newcommand{\npm}{npm\xspace}
\newcommand{\pypi}{PyPI\xspace}
\newcommand{\osv}{OSV\xspace}

\begin{document}

\date{}
\title{\Large \bf Phantom Dependencies: Ghost Busting}

\author{
{\rm Antonio Cyé Angelberg Kambiré}\\
Stanford University
\and
{\rm Julián Rodríguez Cárdenas}\\
Stanford University
\and
{\rm Yasin Yousuf}\\
Stanford University
}

\maketitle

\begin{abstract}
As programming languages and the projects built with them become more complex,
programmers have started to rely on centralized open source package repositories
to provide building blocks for their tasks.  This has led to an intricate
network of dependencies between projects, and that network has repeatedly been
shown to be a source of security risk.  In this work, we investigate
\emph{phantom dependencies}: external packages used in code but not explicitly
acknowledged in a package's manifest.  We study phantom dependencies in the
\npm and \pypi ecosystems using the top 1,000 packages in each repository as
seed sets.  For \npm, we find 118 unique externally managed phantom packages
across 567 seed-closure occurrences after filtering by transitive-closure
membership and maintainer family.  Resolver validation shows that 18 reachable
candidate rows across 5 seeds currently resolve to an affected version, and 15
strict rollback-candidate rows across 6 seeds are reachable in real installs.
For \pypi, a static collector successfully analyzes 801 repositories and finds
593 packages with at least one filtered phantom candidate.  Because the current
\pypi analysis is not version-resolved, its vulnerability results are reported
as advisory-history signals rather than current exploitability.  Overall, our
results show that phantom dependencies are present in both ecosystems, but that
security conclusions require resolver-aware measurement and careful treatment
of static-analysis false positives.
\end{abstract}

\section{Introduction}

Programming projects often rely on various dependencies in order to function.
Generally, projects will have direct dependencies that they use directly, which
in turn use other dependencies, and so on.  The dependencies that the project
eventually depends on are its transitive dependencies.  Phantom dependencies
arise when a project already has a transitive dependency on another package and
later uses that package in its own code without explicitly depending on it.  In
JavaScript, this means that a phantom package is not listed in
\texttt{package.json} but is used through \texttt{require} or \texttt{import}.
In Python, the package is imported in code but is not declared in common
metadata such as \texttt{requirements.txt} or \texttt{pyproject.toml}.

This mismatch matters because manifests are the main place where developers
pin, audit, and update dependencies.  A package that is used but not declared is
harder for the downstream maintainer to control.  A basic question is therefore
the extent to which increasingly dense dependency structures create hidden
dependency relationships and whether those relationships correspond to
security-relevant vulnerability states.

Our study targets that concrete mismatch between dependency information that is
typically tracked and the set of components that a project can actually access
at build and runtime.  In practice, phantom dependencies can remain invisible to
standard dependency declarations and therefore to the workflows developers rely
on for auditing and version management.  We investigate this question in \npm
and \pypi because of their scale and centrality to modern software development.
Concretely, we ask: (1) How prevalent are externally managed phantom
dependencies among heavily used packages in these repositories? (2) Which
phantom dependencies have known vulnerability advisory history, and how often
do current installs resolve to affected versions? (3) How many packages are
susceptible to the rollback attack model in which a transitive dependency
constrains the visible version of a phantom package?

This paper makes three contributions:
\begin{itemize}[leftmargin=*]
  \item A reproducible pipeline for measuring externally managed phantom
  dependency candidates in \npm and \pypi.
  \item A large-scale \npm study of the top 1,000 packages and their dependency
  closure, including explicit depcheck coverage, false-positive filters, and
  resolver-backed validation.
  \item An evaluation of rollback-style exposure that separates advisory-history
  signals, observed affected installs, and strict rollback candidates.
\end{itemize}

\section{Background and Related Work}

Prior work has shown that package ecosystems amplify supply-chain risk through
dense dependency networks.  Zimmermann et al.~\cite{zimmermann2019smallworld}
quantify how installing a single \npm package can extend trust to many packages
and maintainers.  Zerouali et al.~\cite{zerouali2022impact} study vulnerability
propagation in \npm and RubyGems, while Alfadel et al.~\cite{alfadel2023python}
provide a longitudinal analysis of vulnerabilities in Python packages.  Liu et
al.~\cite{liu2022dependencytrees} argue that vulnerability propagation analyses
need repository-specific dependency resolution semantics, a lesson that is
central to our measurement.

Recent practitioner work has also popularized the phrase ``phantom
dependency'' as a source of software bill-of-materials inaccuracy
\cite{endor2023phantom}.  PEP 770 targets a related visibility problem by
standardizing embedded SBOM documents in Python distributions
\cite{pep770}.  Other work explores exploitability along dependency chains:
ChainFuzz generates downstream proof-of-concept exploits for upstream
vulnerabilities~\cite{chainfuzz2025}, and Duress studies build-time supply-chain
hazards in Android resource mediation~\cite{duress2023}.  Our work fits between
these threads: we measure a manifest visibility gap and ask when it translates
into potential security exposure.

\section{Methodology}

\subsection{Problem Formulation: Externally Managed Phantom Dependencies}

For a package $p$ and package name $x$, we call $x$ an externally managed
phantom dependency of $p$ when all of the following hold:
\begin{enumerate}[leftmargin=*]
  \item static analysis reports that $p$ uses $x$;
  \item $x$ is not declared as a direct dependency of $p$;
  \item $x$ is present in $p$'s transitive dependency closure; and
  \item $x$ is not in the same inferred maintainer family as $p$.
\end{enumerate}

The third condition separates phantom dependencies from ordinary missing
dependencies: if the package is not in the closure, it is not plausibly supplied
by the dependency tree.  The fourth condition narrows the analysis from phantom
dependencies in general to externally managed phantom dependencies.  Many
phantom dependencies arise within a single maintainer or organizational
boundary; in such cases, release and security control may remain inside one
family.  We focus on cases where the phantom package is managed outside the
family of the dependent project.

For \npm family inference, we use package scopes such as
\texttt{@babel/*}, GitHub repository owners, and the legacy family metadata in
our database.  For \pypi, the current collector filters obvious local modules,
standard-library imports, build/test tooling, private imports, and ambiguous
namespace roots such as \texttt{google} and \texttt{azure}; it does not yet
construct maintainer families with the same precision as the \npm pipeline.

\subsection{Seed Sets and Data Collection}

Rather than attempting to analyze all packages from \npm and \pypi, we
construct a seed set composed of the top 1,000 most-downloaded packages in each
ecosystem over a six-month window from 2025-05-11 to 2025-11-11.  We focus on
top packages because positive or negative security properties in these packages
can have substantial downstream impact.  For \npm, the checked-in database
contains 12,610 packages reached from the seed set.  We use registry metadata to
build a name-level dependency closure and Dockerized depcheck output~\cite{depcheck}
to identify used-but-missing package names.  The database contains depcheck
results for all 12,610 packages, but only 8,606 runs succeeded; we therefore
treat all prevalence results as lower bounds with explicit coverage fields.

For \pypi, we wrote a fresh static collector that resolves source repositories
from PyPI metadata, clones the repository, parses common manifest files, scans
Python AST imports, and filters standard-library, local, private, build/test,
and artifact-like names.  Of 1,000 seeds, 801 repositories were successfully
analyzed.  The remaining 199 failed because no source repository was identified
or the clone failed.  These failures are important because several high-rank
packages are monorepos or expose metadata URLs that are not direct source
repository URLs.

\subsection{Vulnerability Assessment}

We query \osv~\cite{osv} for vulnerability advisory history.  For \npm, we also
use package versions in the local database and a resolver-backed validation pass
over the security-relevant seeds.  The validation installs each seed with
\texttt{npm install --ignore-scripts}, checks whether the phantom package is
reachable from the package that uses it under Node's module resolution layout,
records the resolved version, and checks that version against OSV affected
ranges.

\subsection{Assessing Susceptibility to Rollback Attack}

We evaluate a rollback-style scenario in which a transitive dependency constrains
the version of a phantom package that a downstream package uses.  We report
three increasingly strict signals:
\begin{description}
  \item[Advisory history.] The phantom package has at least one OSV advisory.
  \item[Affected-range candidate.] A provider dependency range intersects an
  affected range for the phantom package.
  \item[Strict rollback candidate.] A provider range intersects an affected
  advisory range and excludes every known fixed version for that advisory.
\end{description}
For \npm, we further require reachability in an actual install when reporting
our most conservative rollback estimates.  For \pypi, we do not claim rollback
vulnerability because the current pipeline does not resolve installed versions
or provider constraints.

\section{Results}

\subsection{The npm Ecosystem}

We study the top 1,000 most-downloaded \npm packages together with their
transitive dependency closure.  For each package with source information, the
pipeline uses depcheck results to identify packages used by code but missing
from direct dependency declarations.  Raw depcheck output contains several
dominant sources of false positives, including internal packages that are not
published on the registry, packages that do not appear in the dependency
closure, and packages that fall under the same organizational boundary as the
analyzed project.  The results below report the externally managed,
closure-confirmed subset.

\subsubsection{RQ1: npm Prevalence}

\begin{table}[t]
\centering
\small
\begin{tabular}{lrr}
\toprule
Metric & Count & Share \\
\midrule
Seed packages & 1,000 & -- \\
Packages in database & 12,610 & -- \\
Successful depcheck runs & 8,606 & 68.2\% \\
Failed depcheck runs & 4,004 & 31.8\% \\
Seeds with full closure coverage & 585 & 58.5\% \\
Mean seed-closure depcheck coverage & -- & 82.1\% \\
\bottomrule
\end{tabular}
\caption{npm analysis coverage. Failed depcheck runs are counted as unknown,
not as clean packages.}
\label{tab:npm-coverage}
\end{table}

Table~\ref{tab:npm-coverage} shows the main limitation of the \npm analysis.
The database is complete at the metadata level, but depcheck did not succeed for
every package.  As a result, the counts below should be read as lower bounds
over successfully analyzed packages in each seed closure.

\begin{table}[t]
\centering
\small
\begin{tabular}{lr}
\toprule
Metric & Count \\
\midrule
Unique externally managed phantoms & 118 \\
Total seed-closure occurrences & 567 \\
Seeds with any closure-level phantom & 27 \\
Unique root-level phantoms & 107 \\
Root-level phantom occurrences & 233 \\
Seeds with any root-level phantom & 19 \\
Unique phantoms with OSV advisory history & 22 \\
Advisory-history phantom occurrences & 105 \\
\bottomrule
\end{tabular}
\caption{Strict npm phantom dependency results after closure and family
filtering.}
\label{tab:npm-results}
\end{table}

After filtering raw depcheck output by database presence, transitive closure,
and family membership, we find 118 unique externally managed phantom packages
across 567 seed-closure occurrences (Table~\ref{tab:npm-results}).  Only
27/1,000 seeds have at least one closure-level externally managed phantom.

The largest phantom groups are concentrated in familiar JavaScript utility and
visualization packages.  The most frequent names include \texttt{d3-array} (19
occurrences), \texttt{d3-time} (17), \texttt{d3-format} (14),
\texttt{d3-time-format} (14), and \texttt{lodash} (13).  These are not all
security problems; they are evidence that widely used packages sometimes rely on
packages supplied by nearby transitive dependencies.

\subsubsection{RQ2: npm Vulnerability Signals}

Twenty-two unique \npm phantom packages have advisory history in OSV, accounting
for 105 seed-closure occurrences.  The highest-severity set includes
\texttt{lodash}, \texttt{@babel/traverse}, \texttt{d3-color},
\texttt{semver}, \texttt{minimatch}, \texttt{tar}, and \texttt{axios}.  However,
advisory history alone is not equivalent to current exploitability: a package
can have past advisories while the version visible in a current install is
patched.

The resolver validation makes this distinction concrete.  We validated 105
security-relevant \npm occurrence rows across 19 seeds.  All 19 seeds installed
successfully.  Of the 105 rows, 102 were actually reachable from the package
that used the phantom.  Eighteen reachable rows across 5 seeds resolved to a
version affected by at least one OSV advisory.  Thus the most conservative
``currently affected install'' estimate is 5/1,000 seeds in this measurement
snapshot.

\subsubsection{RQ3: npm Rollback Exposure}

\begin{table}[t]
\centering
\small
\begin{tabular}{lrr}
\toprule
Signal & Rows & Seeds \\
\midrule
Affected-range candidates & 177 & 14 \\
Strict range candidates & 26 & 6 \\
Reachable affected-range candidates & 53 & 14 \\
Reachable strict rollback candidates & 15 & 6 \\
Observed affected resolved versions & 18 & 5 \\
\bottomrule
\end{tabular}
\caption{npm rollback-related signals. Rows are package/phantom/provider
relationships; seeds count top-level packages.}
\label{tab:rollback}
\end{table}

Table~\ref{tab:rollback} reports the rollback analysis.  A loose upper bound
finds 14 seeds with a provider range that admits an affected version of a
phantom package.  The stricter advisory-level test finds 26 rows across 6 seeds
where the provider range intersects an affected range and excludes the known
fixed versions for that advisory.  Resolver validation shows that 15 of these
strict rows are reachable in current installs, still across 6 seeds.  The
dominant strict cases involve \texttt{minimatch}, \texttt{lodash},
\texttt{tar}, \texttt{ws}, and \texttt{axios}.

We deliberately describe these as rollback candidates rather than confirmed
exploits.  A full exploit would also require demonstrating that an attacker can
control or compromise the provider dependency and that the downstream code path
uses the vulnerable functionality.  The results nevertheless show that the
rollback model is not merely hypothetical: there are current install layouts in
which a downstream package reaches a phantom through a transitive dependency,
and the provider's declared range can constrain the version available to
downstream phantom imports.

\subsection{The PyPI Ecosystem}

\begin{table}[t]
\centering
\small
\begin{tabular}{lr}
\toprule
Metric & Count \\
\midrule
Seed packages & 1,000 \\
Successful source analyses & 801 \\
Failed/no-repo analyses & 199 \\
Packages with filtered phantom candidates & 593 \\
Unique filtered phantom candidates & 1,657 \\
Total filtered phantom occurrences & 4,688 \\
Mean candidates per successful package & 5.85 \\
Median candidates per successful package & 2 \\
\bottomrule
\end{tabular}
\caption{PyPI static-analysis results after artifact and tooling filters.}
\label{tab:pypi}
\end{table}

The \pypi results are intentionally reported as static-analysis candidates.
After filtering obvious artifacts, 593 of 801 successfully analyzed packages
have at least one phantom candidate (Table~\ref{tab:pypi}).  The most common
candidate names include \texttt{typing-extensions}, \texttt{requests},
\texttt{numpy}, \texttt{packaging}, \texttt{ipython}, \texttt{jinja2},
\texttt{pyyaml}, and \texttt{pydantic-core}.  Many of these are plausible
runtime packages, but some can also appear through optional dependencies,
compatibility code, or packaging-specific import conventions.

OSV advisory-history linkage finds 180 unique \pypi phantom candidates with at
least one advisory, accounting for 1,072 static occurrences across 348 packages.
We do not claim that these packages are currently vulnerable.  The present
\pypi collector does not resolve installed versions, evaluate markers and
extras through pip's resolver, or prove that a particular provider supplies the
imported distribution.  Consequently, the \pypi vulnerability numbers should be
read as prioritization signals for future validation, not as vulnerability
prevalence.

\section{Discussion and Future Work}

Our results characterize the scope of externally managed phantom dependencies
in popular \npm and \pypi packages and quantify several forms of downstream
exposure.  They also show why phantom dependency measurement has to be more
precise than a raw manifest/static-analysis comparison.  A package name that is
reported as missing by a static tool may be absent from the dependency closure,
may be part of the same maintainer family, or may come from build-time and
test-time code that is not relevant to runtime exposure.  Treating these cases
separately is essential for estimating prevalence and security impact.

The \npm analysis is limited by depcheck coverage.  About one third of package
depcheck runs failed, and 415 seeds have at least one failed package in their
closure.  Our \npm prevalence values are therefore lower bounds over the
successfully analyzed portion of each closure.

The \npm database is also name- and latest-version based.  Resolver validation
uses current \npm installs for security-relevant seeds, which improves the
rollback estimate, but the broad closure analysis does not reconstruct every
historical package-lock tree.  Monorepos and source repositories can also differ
from published tarballs.  A future pipeline should analyze exact npm tarballs or
published \texttt{gitHead} commits and run static analysis in the package
subdirectory specified by registry metadata.

For \pypi, the largest limitation is the absence of resolver-backed versions.
Python import names do not always equal distribution names, optional
dependencies and environment markers matter, and many packages use generated
code or compatibility imports.  A version-aware \pypi vulnerability analysis
should use \texttt{pip install --dry-run --report}, map imports to installed
distributions through metadata, and evaluate OSV ranges against resolved
versions.

Several directions follow naturally.  First, a fuller vulnerability assessment
should distinguish build-time, test-time, optional, and runtime code paths and
incorporate additional signals such as maintenance status and package
abandonment.  Second, the rollback attack model should be evaluated with
end-to-end proof-of-concept cases that show how a transitive dependency can
constrain the resolved version of a phantom package in practice.  Third, the
analysis should be scaled beyond the top 1,000 packages in each ecosystem to
measure how phantom dependency behavior changes among less popular packages.
Finally, the results raise an explanatory question: why do projects continue to
use undeclared dependencies even though ecosystem-specific tools can often
detect them?  Understanding that maintenance behavior is necessary for turning
measurement into effective developer guidance.

\section{Conclusion}

In this paper, we provide a large-scale analysis of phantom dependencies in
\npm and \pypi, two of the most widely used package registries for JavaScript
and Python projects.  We find that phantom dependencies are present among
popular packages in both registries, but that their security interpretation
depends strongly on resolver behavior and false-positive filtering.  In \npm, a
strict closure- and family-filtered analysis of the top 1,000 packages finds
118 unique externally managed phantom packages and 567 seed-closure
occurrences.  Resolver validation narrows the version-aware security conclusion
further: 18 reachable rows across 5 seeds currently resolve to affected
versions, and 15 strict rollback-candidate rows across 6 seeds are reachable in
actual installs.  In \pypi, filtered static analysis finds many candidates, but
the security interpretation remains advisory-history only until resolver-backed
version data is added.  Overall, our results suggest that phantom dependencies
are a practical supply-chain visibility problem, and that tooling integration
can improve visibility, control, and remediation when the analysis is
resolver-aware.

\section*{Availability}

The accompanying repository contains the analysis code, generated CSV/JSON
artifacts, and this LaTeX folder.  The main reproducibility commands are
documented in \texttt{README.md}.  The generated artifacts used in this paper
are copied into \texttt{paper\_usenix/data/}.  The \pypi claims use the
\texttt{pypi-collected-*} artifacts.

\bibliographystyle{plain}
\bibliography{references}

\end{document}
